\section{Reproduction}

\subsection{Description of the method}
We reimplemented the GVP-WM planning pipeline around the public DINO-WM
checkpoint. A video plan is first encoded into latent states. The planner then
jointly optimizes latent states and bounded action variables, while an
augmented Lagrangian penalty encourages each optimized transition to match the
action-conditioned DINO-WM dynamics. The optimized actions are executed in the
environment in a receding-horizon manner.

We evaluated three sources of video guidance. ORACLE uses videos obtained from
expert environment rollouts. WAN-0S uses zero-shot videos generated by Wan2.1
FLF2V from the start and target frames. WAN-FT uses our available fine-tuned Wan
checkpoints. In all cases, generated frames are cropped back to the environment
observation format before being encoded by DINO-WM. We kept the original success
criteria and counted video-generation or decoding failures as failed trials.

\subsection{Experiments}
We ran the Push-T and Wall settings corresponding to Table 2 of the original
paper. Each entry below is a success rate over 50 sampled initial-goal pairs.
For Push-T we evaluate horizons \(T\in\{25,50,80\}\); for Wall we evaluate
\(T\in\{25,50\}\). The table reports our result followed by the paper result in
parentheses.

\begin{table}[h]
\centering
\small
\setlength{\tabcolsep}{4pt}
\begin{tabular}{lccccc}
\toprule
Method & Push-T 25 & Push-T 50 & Push-T 80 & Wall 25 & Wall 50 \\
\midrule
ORACLE & .68 (.98) & .20 (.72) & .06 (.36) & .94 (1.00) & .94 (1.00) \\
WAN-0S & .48 (.56) & .02 (.12) & .00 (.04) & .84 (.86) & .86 (.76) \\
WAN-FT & .52 (.80) & .02 (.30) & .00 (.06) & .96 (.94) & .92 (.90) \\
\bottomrule
\end{tabular}
\caption{Reproduction results. Each cell shows our success rate with the
reported paper value in parentheses.}
\label{tab:reproduction-results}
\end{table}

\subsection{Discussion and analysis}
The reproduction is partially successful but not paper-faithful on all tasks.
Wall is close to the reported performance: both ORACLE and WAN-FT reach above
\(0.9\), and WAN-0S is near or above the reported Wall numbers. This suggests
that the environment interface, frame handling, and video-to-planning pipeline
are broadly functional.

Push-T remains the main gap. Even ORACLE guidance is far below the paper
upper bound, especially for longer horizons. Since ORACLE removes video
generation quality from the loop, the failure is unlikely to be explained only
by Wan video quality. The more likely bottleneck is the action recovery and
optimization protocol for contact-rich Push-T trajectories, including the
interaction between action scaling, action regularization, latent/proprio
handling, and ALM optimization. For this reason, we do not claim a full
reproduction of Table 2; instead, we use the reproduced pipeline as the baseline
for the feasibility-prior experiments in the next section.
